%% kapa: exam=1 round=25 year=2014 period=1 form=A subject=civil_law subject_ko=민법 seq=4 number=4
%% source: Q-Net 감정평가사 기출문제 (cst00309) artlSeq=5101207 file=extracted/1008/5101207/2014년 제25회 감정평가사 1차시험 기출문제/1교시(A형).hwp
%% conversion: kordoc 4.24.51 -> tools/gen1.py (auto); review=proofread-sonnet; flags=none
\begin{question}{4}
비법인사단에 관한 설명으로 옳지 않은 것은? (다툼이 있으면 판례에 의함)
\begin{choices}
  \item 비법인사단이 타인 간의 금전채무를 보증하는 행위는 총유물의 관리․처분행위라고 볼 수 없다.
  \item 사단법인의 하부조직이 스스로 단체로서의 실체를 갖추고 독자적인 활동을 하고 있다면, 사단법인과는 별개의 독립된 비법인사단으로 볼 수 있다.
  \item 비법인사단의 구성원은 총유재산에 대한 보존행위로서 각자 소를 제기할 수 있다.
  \item 비법인사단이 해산하였더라도 청산사무가 완료될 때까지는 청산의 목적범위 내에서 권리․의무의 주체가 된다.
  \item 사단법인에 관한 민법규정 중 법인격을 전제로 한 규정을 제외하고는 비법인사단에도 유추적용된다.
\end{choices}
\end{question}
\answer{3}
